% sec-zeros.tex -- the certified zeros of A083837 and
% A083838. Standalone fragment, \input by
% A083833-UNIFIE-paper.tex.
% Environments and notation inherited from sec-carry.tex; the
% guards below are no-ops once defined there.
\providecommand{\revW}{\mathrm{rev}}
\makeatletter
\@ifundefined{theorem}{\newtheorem{theorem}{Theorem}}{}
\@ifundefined{lemma}{\newtheorem{lemma}[theorem]{Lemma}}{}
\@ifundefined{corollary}{\newtheorem{corollary}[theorem]{Corollary}}{}
\@ifundefined{remark}{\newtheorem{remark}[theorem]{Remark}}{}
\makeatother

\section{The certified zeros of A083837 and A083838}
\label{sec:zeros}

The last two entries of the family scan the multiplier.
A083837~\cite{oeis} (Murthy and Srikanth, May 2003;
corrected and extended by R.~Chandler) is quoted in
Section~\ref{sec:intro}: $a(k)$ is the smallest
palindrome $p$ such that $kp+1$ is palindromic, or $0$
if none exists. Its $104$ displayed terms vanish at the
$21$ indices
\[
 k \in Z = \{12,17,23,31,37,41,42,48,51,53,61,64,71,
 74,79,81,86,91,96,101,103\},
\]
where Chandler's comment reports only the absence of a
palindrome below $10^{15}$. A displayed zero is thus
finite evidence, not nonexistence; a later witness could
refute it. This section upgrades each of the $21$ zeros
to a universal theorem, by an obstruction on the two
ends of the candidate palindrome.

\paragraph{The prefix-suffix obstruction.}
Let $p$ be an $m$-digit palindrome, $1 \le w \le m$, and
let $P$ be its $w$-prefix (the integer formed by its
first $w$ digits). Write $\revW_w(n)$ for the integer
read from the last $w$ digits of $n$ in reverse order.
Assume $q = kp+1$ is a palindrome of $m'$ digits. Then:
\begin{itemize}
\item the prefix bounds the value:
$kP\,10^{m-w} \le q < k(P+1)\,10^{m-w}$;
\item the last $w$ digits of $p$ form $\revW_w(P)$
(palindromicity of $p$), i.e.\
$p \equiv \revW_w(P) \pmod{10^w}$, so the last $w$
digits of $q$ are those of $k\,\revW_w(P) + 1$; as $q$ is
itself palindromic of length $m' \ge m \ge w$, its
$w$-prefix is the reversal of its last $w$ digits:
$L = \revW_w\bigl(k\,\revW_w(P) + 1\bigr)$. If
$L < 10^{w-1}$ the last digit of $q$, hence its leading
digit, would be $0$ -- impossible, so $L \ge 10^{w-1}$
and $L\,10^{m'-w} \le q < (L+1)\,10^{m'-w}$;
\item writing $s = m' - m \ge 0$, one has $s \le S$ for
any $S$ with $k < 10^S$, since
$q < k\,10^m + 1 \le 10^{m+S}$; crossing the two windows
($L\,10^{m'-w} \le q < k(P+1)\,10^{m-w}$ and
$kP\,10^{m-w} \le q < (L+1)\,10^{m'-w}$) and dividing by
$10^{m-w}$ then forces
\begin{equation}
 L\,10^{s} \;\le\; k\,(P+1)
 \qquad\text{and}\qquad
 k\,P \;<\; (L+1)\,10^{s} .
 \label{eq:overlap}
\end{equation}
\end{itemize}

\begin{lemma}[border obstruction]
\label{lem:border}
If, for a candidate $w$-prefix $P$, either
$L < 10^{w-1}$ (the required $w$-prefix of $q$ would
start with $0$, i.e.\ $q$ would end in $0$) or no
$s \in \{0,\dots,S\}$ satisfies \eqref{eq:overlap},
then no $m$-digit palindrome $p$ with $w$-prefix $P$
can have $kp+1$ palindromic. In particular, if every
$w$-prefix $P \in \{10^{w-1},\dots,10^{w}-1\}$ is
obstructed, no palindrome $p$ of length $\ge w$ has
$kp+1$ palindromic.
\end{lemma}

The lemma prunes the prefix tree level by level: the
survivors of depth $w$ seed the candidate $(w+1)$-prefixes
$10P + d$ for $d \in \{0,\dots,9\}$. The residual cases
$m < w$ are finite: an $m$-digit palindrome is its own
$m$-prefix, so it must appear among the level-$m$
survivors, be palindromic as a prefix, and pass the
direct check on $kp+1$. Table~\ref{tab:zeros} gives the
certificate tree for each of the $21$ indices: the depth
$w$ at which the survivor list becomes empty, and the
survivors $E_t$ at each shallower level.

\begin{table}[ht]
\centering
\small
{\footnotesize\setlength{\tabcolsep}{4pt}
\begin{tabular}{>{\raggedright\arraybackslash}p{4.4cm}|c|%
>{\raggedright\arraybackslash}p{10.2cm}}
\hline
$k$ & depth & survivors \\
\hline
$12$, $31$, $37$, $41$, $42$, $51$, $53$, $61$, $64$,
$71$, $81$, $86$, $91$, $96$ & $1$ & none: every leading
digit $a_0 \in \{1,\dots,9\}$ is obstructed \\
$17$ & $2$ & $E_1 = \{2\}$; $m = 1$: $17\cdot2+1 = 35$
not palindromic; $E_2 = \varnothing$ \\
$23$ & $2$ & $E_1 = \{1\}$; $m = 1$: $23\cdot1+1 = 24$
not palindromic; $E_2 = \varnothing$ \\
$74$ & $2$ & $E_1 = \{6,9\}$; $m = 1$: $445$, $667$
not palindromic; $E_2 = \varnothing$ \\
$101$ & $2$ & $E_1 = \{1,\dots,8\}$ ($S = 3$);
$m = 1$ excluded case by case; $E_2 = \varnothing$ \\
$103$ & $2$ & $E_1 = \{5\}$ ($S = 3$); $m = 1$:
$516$ not palindromic; $E_2 = \varnothing$ \\
$48$ & $3$ & $E_1 = \{1\}$; $m = 1$: $48\cdot1+1 = 49$
not palindromic; $E_2 = \{19\}$; $m = 2$: candidate
prefix $19$ not a palindrome ($9 \ne 1$);
$E_3 = \varnothing$ over $\{190,\dots,199\}$ \\
$79$ & $3$ & $E_1 = \{6\}$; $m = 1$: $79\cdot6+1 = 475$
not palindromic; $E_2 = \{65\}$; $m = 2$: candidate
prefix $65$ not a palindrome; $E_3 = \varnothing$
over $\{650,\dots,659\}$ \\
\hline
\end{tabular}}
\caption{Certificate trees of the $21$ certified zeros
of A083837. $E_t$ is the list of $t$-prefixes escaping
the obstruction at level $t$; $m < w$ cases are
finished by direct products.}
\label{tab:zeros}
\end{table}

\begin{theorem}[the $21$ zeros of A083837]
\label{thm:zeros}
For every $k \in Z$ there is no palindrome $p$ of any
length such that $kp+1$ is a palindrome.
\end{theorem}

\begin{proof}
For each $k$, the obstruction of
Lemma~\ref{lem:border} is evaluated on the candidate
prefixes level by level: at depth $1$ the list
$E_1 = \varnothing$ for the fourteen indices of the
first row of Table~\ref{tab:zeros}; at depth $2$ for
$k \in \{17,23,74,101,103\}$ the level-$1$ survivors
shown are excluded at length $m = 1$ by the displayed
products and the next level is empty; at depth $3$ for
$k \in \{48,79\}$ the chains shown close the tree.
Every case is kernel-checked in \texttt{A083837.lean}
(Table~\ref{tab:lean}).
\end{proof}

\begin{corollary}[the $21$ zeros of A083838]
\label{cor:zeros838}
For every $k \in Z$ there is no pair of palindromes
$(p,q)$ with $q = kp+1$: the image entry A083838
vanishes at the same indices.
\end{corollary}

The companion entry A083838 displays $66$ terms; its
$12$ shown zeros, at indices $12$, $17$, $23$, $31$,
$37$, $41$, $42$, $48$, $51$, $53$, $61$ and $64$, are
exactly the displayed part of the certified set, while
the remaining nine indices $71$, $74$, $79$, $81$,
$86$, $91$, $96$, $101$ and $103$ lie beyond its
displayed window and are settled here on the same
footing.

\begin{proof}
Each corollary is the direct negation of the existence
of a member pair, a one-line consequence of
Theorem~\ref{thm:zeros}; the $21$ statements are
kernel-checked in \texttt{A083838.lean}.
\end{proof}

\begin{remark}[scope of the certification]
\label{rem:scopezeros}
Theorem~\ref{thm:zeros} is a universal nonexistence
statement: it rules out every palindrome of every
length, where the displayed zero reported only
$p < 10^{15}$. It says nothing about the nonzero
displayed values of A083837 or A083838, which are not
certified here; and it does not address multipliers
$k > 104$, beyond the displayed range of the entry.
Consistently, an independent enumeration found no
witness palindrome below $10^{11}$ for any of the $21$
indices -- agreement with the finite-search comment,
though never used as evidence.
\end{remark}
